\scanpage{49}

Soit

\[
\unsure{\mathcal X'|\mathfrak X'}:\mathrm{Sch}_{/X}\longrightarrow\mathrm{Ens}\qquad\text{dont la val. sur }X\text{ soit}
\]

\[
\mathcal X'(X)=\left\{x\in\Gamma\underline L\ \middle|\ \begin{aligned}x^2&=\delta\\x|X_0&=T_0\end{aligned}\right\}\qquad\begin{gathered}\text{et idem via}\\Y\longrightarrow X.\end{gathered}
\]

\[
\deleted{\unsure{\text{Soit}}}\quad\mu'_{2\unsure{X|x}}\subset\mu_{2X}:\mathrm{Sch}_{/X}\longrightarrow\text{Groupes comm}
\]

\[
\mu'_2\unsure{=}(X)=\{\lambda\in\Gamma\underline{\mathscr O}_X^*\mid\lambda^2=1,\ \lambda_0=1\}
\]

% slanted
\margin{\uline{NB} $\mathcal X'$ et $\mu'_2$ sont des faisceaux fpqc, \unsure{$\varnothing$|$\mathscr O$} \uline{pas} repr. \deleted{\unsure{si}} $\exists\,x\in X$ tel que, \deleted{\unsure{soit}} \deleted{\unsure{nul soit}} $2$ soit ni inv. ni nul ds $\mathscr O_{X,x}$, i.e. $2\in\mathfrak m_{X,x}$, $2\ne0$}

Alors $\mu'_2$ opère sur $\mathcal X'$, et si $\delta$ est inv. i.e. si $X'/X$ étale, alors $\mathcal X'$ est un torseur sous $\mu'_2$ (alors que $X'/X$ est un torseur sous $(\mathbb Z/2\mathbb Z)_X$).

On a compatibilité

\diagram{\begin{tikzcd}[ampersand replacement=\&]
X' \arrow[r,"f"] \& \mathcal X' \\
(\mathbb Z/2\mathbb Z)_X \arrow[r] \& \mu'_{2X}
\end{tikzcd}}

i.e.

\[
f(\bar x)=-\deleted{\unsure{\bar x}}f(x)
\]

où

\[
z\longmapsto-z:\quad\mathcal X'\longrightarrow\mathcal X'
\]

est l'\deleted{application}\unsure{automorphisme|opération} induit par la section $-1$ de $\mu'_2$ sur $X$ \deleted{\illegible{2}}.

\smallskip\centerline{\rule{1em}{0.3pt}}\smallskip

Supposons $2$ $\underline{\mathscr O}_X$-régulier. Alors la donnée des rev. quadratiques $X'/X$ \uline{équivaut} à celle des

\[
\left\{\begin{aligned}
&\underset{\text{inv.}}{L},\quad\delta\in\Gamma L^{\otimes 2},\quad T_0\in\Gamma(L_0)\qquad(\text{où }L_0=L|X_0)\\
&\text{satisfaisant}\qquad T_0^2\equiv\delta\quad(4)
\end{aligned}\right.
\]

On veut décrire les sections $x$ de $X'/X$, en termes de la section $z$ de $L$, dans le cas où $X'\longrightarrow\mathcal X'$ pas iso i.e. où $X_0\ne\varnothing$. \deleted{On aura} \inserted{Bien sûr,} on aura

% archive sheet number: 48
